\documentclass[11pt,a4paper]{article}
\usepackage[margin=2.5cm]{geometry}
\usepackage{amsmath,amsthm,amssymb}
\usepackage{graphicx}
\usepackage{booktabs}
\usepackage[hidelinks]{hyperref}
\usepackage{cleveref}
\usepackage[backend=biber,style=numeric,urldate=long]{biblatex}
\addbibresource{references.bib}
\graphicspath{{figures/}}
\setlength{\emergencystretch}{3em}
\setcounter{biburlnumpenalty}{100}\setcounter{biburlucpenalty}{100}\setcounter{biburllcpenalty}{100}

\theoremstyle{definition}
\newtheorem{definition}{Definition}[section]
\newtheorem{example}[definition]{Example}
\theoremstyle{plain}
\newtheorem{theorem}[definition]{Theorem}
\newtheorem{lemma}[definition]{Lemma}
\theoremstyle{remark}
\newtheorem{remark}[definition]{Remark}
\theoremstyle{plain}
\newtheorem{proposition}[definition]{Proposition}

\title{Adjacency matrices, walks, paths, and connectivity: study notes}
\author{}
\date{November 2026}

\begin{document}
\maketitle

\begin{abstract}
For a finite simple undirected graph with $N$ vertices and adjacency matrix $A$, we answer three questions. The entry $(A^k)_{ij}$ is the number of walks of length $k$ from vertex $n_i$ to vertex $n_j$. The graph is connected exactly when every entry off the diagonal of $A + A^2 + \dots + A^{N-1}$ is positive. A path between $n_i$ and $n_j$ is a walk between them in which no vertex occurs twice. Both results are proved in full and worked on two small graphs.
\end{abstract}

\section{Introduction}\label{sec:introduction}

Let a finite simple undirected graph have $N$ vertices and be given by its adjacency matrix $A$, a table of zeros and ones that records which vertices are joined by an edge. These notes answer three questions about it. (1) What does the entry $(A^k)_{ij}$ of the $k$-th power of $A$ represent? (2) How can one tell from $A$ alone whether the graph is connected? (3) What is a path between two vertices $n_i$ and $n_j$?

The answers are as follows. (1) The entry $(A^k)_{ij}$ is the number of walks of length $k$ from $n_i$ to $n_j$ (\Cref{thm:walks}); it counts walks, not paths, as \Cref{ex:board} shows on a three-vertex graph. (2) The graph is connected exactly when every entry off the diagonal of $A + A^2 + \dots + A^{N-1}$ is positive (\Cref{thm:test}). (3) A path between $n_i$ and $n_j$ is a walk from $n_i$ to $n_j$ in which no vertex occurs twice (\Cref{def:path}).

\Cref{sec:definitions} fixes the notation and the definitions, in the order of the course textbook \cite{petrosyan_graph_theory}; textbook pages are printed pages, one less than the PDF page. \Cref{sec:walks} proves the answer to question 1, \Cref{sec:connectivity} the answer to question 2, \Cref{sec:examples} works both on two graphs, and \Cref{sec:limits} states what the test does not give. The textbook's own word ``path'' has a different meaning from the one used here; \Cref{rem:trap} says which.

\section{Setting and definitions}\label{sec:definitions}

\begin{definition}[Graph and adjacency; textbook Definitions 1.1.1 and 1.1.4]\label{def:graph}
A \emph{graph} is a pair $G = (V, E)$ of a non-empty finite set $V$ of \emph{vertices} and a set $E$ of \emph{edges}; it is \emph{simple undirected} when each edge is an unordered pair of two distinct vertices, so there are no loops and no repeated edges \cite[Definition 1.1.1, p.~9; Definition 1.1.4, p.~11]{petrosyan_graph_theory} \cite[Definition 1.1.1]{galway_adjacency_notes}. All graphs in this paper are finite, simple, and undirected. Write $V = \{n_1, \dots, n_N\}$ with $N \ge 1$ (the textbook writes $v_i$), and write $n_i n_j = n_j n_i$ for an edge. Two vertices are \emph{adjacent} if they form an edge.
\end{definition}

\begin{definition}[Adjacency matrix; textbook, p.~11, no number of its own]\label{def:matrix}
The \emph{adjacency matrix} $A$ of $G$ is the $N \times N$ matrix with $A_{ij} = 1$ if $n_i$ and $n_j$ are adjacent and $A_{ij} = 0$ otherwise \cite[p.~11]{petrosyan_graph_theory}. Edges are unordered pairs of distinct vertices, so $A_{ij} = A_{ji}$ and $A_{ii} = 0$ \cite[p.~11]{petrosyan_graph_theory}.
\end{definition}

For $N \times N$ matrices $X, Y$ the product has entries $(XY)_{ij} = \sum_{r=1}^{N} X_{ir} Y_{rj}$. The powers of $A$ are $A^0 = I$, the identity matrix (ones on the diagonal, zeros elsewhere), and $A^{k+1} = A^k A$. The entry in row $i$ and column $j$ of $A^k$ is written $(A^k)_{ij}$.

\begin{definition}[Walk and length; textbook Definition 1.2.10]\label{def:walk}
A \emph{walk} of \emph{length} $k \ge 0$ from $n_i$ to $n_j$ is a sequence $(v_0, v_1, \dots, v_k)$ of vertices with $v_0 = n_i$, $v_k = n_j$, and $v_{s-1}$ adjacent to $v_s$ for $1 \le s \le k$. The length $k$ is the number of edges counted with repetition.
\end{definition}

The textbook writes a walk as an alternating sequence of vertices and edges $u_0$, $e_1$, $u_1$, \dots, $e_k$, $u_k$, with $e_s = u_{s-1}u_s$ \cite[Definition 1.2.10, p.~21]{petrosyan_graph_theory}. In a simple graph two vertices carry at most one edge, so the edges are determined by the vertices and the two forms describe the same walks (a derived remark); Grinberg also writes walks as vertex sequences \cite[Definition 2.9.1(a), p.~36]{grinberg_graph_theory}. In a walk both vertices and edges may repeat \cite[p.~22]{petrosyan_graph_theory}. This paper allows $k = 0$: the single vertex $(n_i)$ is a walk of length $0$ from $n_i$ to itself. That is the paper's convention; it agrees with the textbook's property $d(u,v) = 0$ exactly when $u = v$ \cite[p.~23]{petrosyan_graph_theory} and with Grinberg's definition \cite[Example 2.9.2, p.~37]{grinberg_graph_theory}. Since $A$ is symmetric, a walk read backwards is a walk.

\begin{definition}[Path; textbook Definition 1.2.13, ``simple path'']\label{def:path}
A \emph{path} is a walk whose vertices are pairwise distinct. A \emph{path between $n_i$ and $n_j$}, written from $n_i$ to $n_j$, is a path $(v_0, \dots, v_k)$ with $v_0 = n_i$ and $v_k = n_j$ \cite[Definition 1.2.13, p.~22]{petrosyan_graph_theory} \cite[Definition 2.9.1(c),(d)]{grinberg_graph_theory} \cite{bai_bimsa_adjacency}. This is the meaning of ``path'' throughout the paper.
\end{definition}

\begin{remark}\label{rem:trap}
The textbook's own word ``path'' means something weaker: a walk whose edges are pairwise distinct \cite[Definition 1.2.12, p.~22]{petrosyan_graph_theory}, called a \emph{trail} in English. This paper never uses ``path'' in that sense; it writes ``trail''. By \Cref{def:path} each vertex $n_i$ has exactly one path to itself, $(n_i)$, of length $0$. Reversing a path from $n_i$ to $n_j$ gives a path from $n_j$ to $n_i$, so ``between'' is symmetric.
\end{remark}

\begin{definition}[Connected; textbook Definition 2.1.1, restated]\label{def:connected}
$G$ is \emph{connected} if for every pair $n_i, n_j$ of vertices there is a path from $n_i$ to $n_j$ \cite[Definition 2.1.1, p.~48]{petrosyan_graph_theory}.
\end{definition}

The textbook states this with its trail in place of the path; \Cref{prop:same} shows that the two statements describe the same graphs. For $i = j$ the length-$0$ path settles the condition, so only pairs with $i \ne j$ carry content.

\section{Walk counting: the answer to question 1}\label{sec:walks}

\begin{theorem}[Walk counting]\label{thm:walks}
Let $W_k(i,j)$ be the number of walks of length $k$ from $n_i$ to $n_j$. For every integer $k \ge 0$ and all $i, j$ in $\{1, \dots, N\}$, $(A^k)_{ij} = W_k(i,j)$.
\end{theorem}

Galway, Gotti, and Bai state this result \cite[Theorem 1.1.3, p.~4]{galway_adjacency_notes} \cite[Proposition 2, p.~1]{gotti_18211_lecture27} \cite[Theorem 1, p.~1]{bai_bimsa_adjacency}, and the proof below is this paper's own; Galway's version is for $k \ge 1$.

\begin{proof}
By induction on $k$.

\emph{Base case} $k = 0$. A walk of length $0$ from $n_i$ to $n_j$ is a single vertex $v_0$ with $v_0 = n_i$ and $v_0 = n_j$. It exists, and is then unique, exactly when $i = j$. So $W_0(i,j) = 1$ if $i = j$ and $0$ otherwise, which is the $(i,j)$ entry of $I = A^0$.

\emph{Induction hypothesis.} For some $k \ge 0$, $(A^k)_{ir} = W_k(i,r)$ for all $i, r$.

\emph{Step.} Fix $i, j$. A walk $(v_0, \dots, v_{k+1})$ of length $k+1$ from $n_i$ to $n_j$ has exactly one last-but-one vertex $v_k = n_r$, for one $r$ in $1, \dots, N$. Sort the walks by this $r$ and fix $r$. Dropping the last vertex gives a walk $(v_0, \dots, v_k)$ of length $k$ from $n_i$ to $n_r$, and the original walk is recovered by appending $n_j$. Conversely, appending $n_j$ to a walk of length $k$ from $n_i$ to $n_r$ gives a walk exactly when $n_r$ and $n_j$ are adjacent, that is when $A_{rj} = 1$. So the number of walks of length $k+1$ with last-but-one vertex $n_r$ is $W_k(i,r)$ if $A_{rj} = 1$ and $0$ if $A_{rj} = 0$; in both cases it is $W_k(i,r)\,A_{rj}$. Every walk has exactly one $r$, so no walk is missed or counted twice, and
\[
W_{k+1}(i,j) = \sum_{r=1}^{N} W_k(i,r)\,A_{rj} = \sum_{r=1}^{N} (A^k)_{ir} A_{rj} = (A^{k+1})_{ij},
\]
by the hypothesis and the definition of the product. This completes the induction.
\end{proof}

Nothing in the proof asks the vertices of the walk to be distinct, which is why the entries count walks and not paths.

\begin{example}[The board graph]\label{ex:board}
Let $G_1$ have $N = 3$ vertices and the edges $n_1n_2$ and $n_2n_3$. Its adjacency matrix and the square of it are
\[
A = \begin{pmatrix} 0&1&0\\ 1&0&1\\ 0&1&0 \end{pmatrix}, \qquad
A^2 = \begin{pmatrix} 1&0&1\\ 0&2&0\\ 1&0&1 \end{pmatrix}
\]
(computed by script, which also checks them against a brute-force count of walks \cite{m1_walks_script}). By \Cref{thm:walks} the entries of $A^2$ count walks of length $2$.
$(A^2)_{11} = 1$ counts the single walk $n_1, n_2, n_1$. It repeats $n_1$, so it is not a path; the only path from $n_1$ to $n_1$ is $(n_1)$, of length $0$. So this entry counts a walk that is not a path.
$(A^2)_{22} = 2$ counts $n_2, n_1, n_2$ and $n_2, n_3, n_2$, neither a path.
$(A^2)_{13} = 1$ counts $n_1, n_2, n_3$, a path.
Hence the statement ``$(A^k)_{ij}$ is the number of paths of length $k$'' is wrong for this graph.
\end{example}

\section{A connectivity test from $A$: the answer to question 2}\label{sec:connectivity}

A positive $(A^k)_{ij}$ says by \Cref{thm:walks} only that a walk of length $k$ exists, and a walk may revisit vertices, while \Cref{def:connected} asks for paths. The next lemma closes that gap and also bounds the lengths to be checked.

\begin{lemma}[Walk contains path]\label{lem:walkpath}
Let $G$ have $N$ vertices.
(a) If there is a walk of length $k$ from $n_i$ to $n_j$ ($i = j$ allowed), there is a path from $n_i$ to $n_j$ of length at most $k$ whose vertices all occur in the walk.
(b) A path in $G$ has length at most $N - 1$.
\end{lemma}

Grinberg proves (a) and states (b) \cite[Corollary 2.9.7, p.~39; p.~40]{grinberg_graph_theory}; the textbook's Lemma 1.2.1 is stated for $u \ne v$ and has no length bound \cite[p.~23]{petrosyan_graph_theory}. The proofs below are this paper's own.

\begin{proof}
(a) By strong induction on $k$. \emph{Base case} $k = 0$: the walk is $(n_i)$, so $n_j = n_i$, and $(n_i)$ has no repeated vertex, so it is a path of length $0$. \emph{Step:} let $k \ge 1$ and assume (a) for every walk of smaller length. Let $w = (w_0, \dots, w_k)$ be a walk from $n_i$ to $n_j$. If no vertex repeats in $w$, then $w$ is a path of length $k$. Otherwise $w_a = w_b$ for some $a < b$. Form $w' = (w_0, \dots, w_a, w_{b+1}, \dots, w_k)$, which is $(w_0, \dots, w_a)$ when $b = k$. It is a walk: consecutive terms inside the two kept blocks were adjacent in $w$, and the one new consecutive pair is $(w_a, w_{b+1}) = (w_b, w_{b+1})$, adjacent in $w$. It starts at $w_0 = n_i$, and it ends at $w_k = n_j$ (when $b = k$ it ends at $w_a = w_b = w_k$). Its length is $k - (b - a) < k$. By the induction hypothesis $w'$ contains a path from $n_i$ to $n_j$ of length at most $k - (b-a) \le k$, with vertices occurring in $w'$ and so in $w$.

(b) A path of length $k$ has $k + 1$ terms $w_0, \dots, w_k$, pairwise distinct vertices of $G$. $G$ has $N$ vertices, so $k + 1 \le N$.
\end{proof}

\begin{proposition}[The textbook's connected and this paper's agree]\label{prop:same}
For all $n_i, n_j$ in a simple graph, there is a path from $n_i$ to $n_j$ (\Cref{def:path}) if and only if there is a trail from $n_i$ to $n_j$. Hence the graphs that are connected by \Cref{def:connected} are the graphs that are connected by the textbook's Definition 2.1.1, which uses its own ``path'', the trail of \Cref{rem:trap}.
\end{proposition}

\begin{proof}
(Derived here.) \emph{Path implies trail.} Let $(v_0, \dots, v_k)$ be a path and suppose two of its edges are equal, $\{v_{s-1}, v_s\} = \{v_{t-1}, v_t\}$ with $s < t$. Then $v_{s-1}$ is $v_{t-1}$ or $v_t$. But $s - 1 < t - 1 < t$, so $v_{s-1}$ sits at a different position from both, and the vertices of a path are distinct: a contradiction. So the edges are pairwise distinct and the path is a trail. \emph{Trail implies path.} A trail from $n_i$ to $n_j$ is a walk from $n_i$ to $n_j$, so by \Cref{lem:walkpath}(a) it contains a path from $n_i$ to $n_j$. The two directions give the equivalence for each pair, hence for the quantified definitions.
\end{proof}

\begin{theorem}[Sum test]\label{thm:test}
Let $N \ge 2$ and $S = A + A^2 + \dots + A^{N-1}$. The graph is connected if and only if $S_{ij} > 0$ for all $i \ne j$.
\end{theorem}

Paramadevan and Sotheeswaran state this test; their proof moves from a non-zero entry to a path without the step of \Cref{lem:walkpath}(a) \cite[Theorem 2, p.~17]{paramadevan2021adjacency}. The proof below supplies that step and is this paper's own.

\begin{proof}
\emph{Step 0.} By \Cref{thm:walks} each $(A^k)_{ij}$ is a number of walks, a non-negative integer. So $S_{ij} = W_1(i,j) + \dots + W_{N-1}(i,j)$, and $S_{ij} > 0$ if and only if for some $k$ with $1 \le k \le N-1$ there is a walk of length $k$ from $n_i$ to $n_j$.

\emph{Connected implies $S_{ij} > 0$.} Let $i \ne j$. By \Cref{def:connected} there is a path $P$ from $n_i$ to $n_j$, of length $k$ say. Since $n_i \ne n_j$, $P$ has at least two vertices, so $k \ge 1$, and $k \le N - 1$ by \Cref{lem:walkpath}(b). $P$ is a walk of length $k$, so $S_{ij} \ge W_k(i,j) \ge 1$.

\emph{$S_{ij} > 0$ for all $i \ne j$ implies connected.} Take any pair $n_i, n_j$. If $i = j$, the path $(n_i)$ of length $0$ serves. If $i \ne j$, Step 0 gives a walk from $n_i$ to $n_j$, and \Cref{lem:walkpath}(a) gives a path from $n_i$ to $n_j$. Every pair has a path, so $G$ is connected.
\end{proof}

The bound $N-1$ is used only in the first direction. The diagonal is left out because it is not needed (each vertex has its length-$0$ path) and because it can be zero in a connected graph: for $N = 2$ and one edge, $S = A$ has zero diagonal. For $N = 1$ the graph is connected and no pair has $i \ne j$, so the condition holds vacuously, while $S$ is an empty sum; the test is therefore stated for $N \ge 2$.

\begin{remark}[The form $(I+A)^{N-1}$]\label{rem:binom}
Gotti and Bai state the test with $R = (I + A)^{N-1}$ in place of $S$ \cite[Proposition 3, p.~2]{gotti_18211_lecture27} \cite[Theorem 2, p.~1]{bai_bimsa_adjacency}. The following is derived here. Since $I$ and $A$ commute, the binomial theorem gives $R = \sum_{k=0}^{N-1} \binom{N-1}{k} A^k$. Every coefficient is at least $1$ and every $A^k$ has non-negative entries, so for $i \ne j$ the entry $R_{ij}$ is positive exactly when $S_{ij}$ is (the term $k = 0$ adds only to the diagonal), and each $R_{ii} \ge 1$. So $R$ has no zero entry exactly when the test of \Cref{thm:test} holds. By script, $R$ has rows $(2,2,1), (2,3,2), (1,2,2)$ for $G_1$ and $(2,2,0), (2,2,0), (0,0,1)$ for $G_2$ of \Cref{sec:examples}.
\end{remark}

\section{Worked examples}\label{sec:examples}

All matrices below are computed by script, which also checks each $S_{ij}$ against a brute-force count of walks and each answer (connected or not) against a brute-force search for paths; the Notes on sources say how to rerun it.

\begin{example}[The test on the example graph $G_1$]\label{ex:test1}
$G_1$ is the graph of \Cref{ex:board}, $N = 3$, so $S = A + A^2$:
\[
A = \begin{pmatrix} 0&1&0\\ 1&0&1\\ 0&1&0 \end{pmatrix},\quad
A^2 = \begin{pmatrix} 1&0&1\\ 0&2&0\\ 1&0&1 \end{pmatrix},\quad
S = \begin{pmatrix} 1&1&1\\ 1&2&1\\ 1&1&1 \end{pmatrix}.
\]
The off-diagonal entries are $S_{12} = 1 + 0$ (the walk $n_1, n_2$), $S_{13} = 0 + 1$ (the walk $n_1, n_2, n_3$), and $S_{23} = 1 + 0$ (the walk $n_2, n_3$); $S$ is symmetric, so $S_{21}, S_{31}, S_{32}$ equal them. All six are positive, so \Cref{thm:test} says $G_1$ is connected, as the matrix suggests: $n_2$ is adjacent to $n_1$ and to $n_3$. The diagonal is not read: $S_{11} = 1$ and $S_{22} = 2$ play no role.
\end{example}

\begin{example}[The test on a disconnected graph $G_2$]\label{ex:test2}
Let $G_2$ have $N = 3$ vertices and the single edge $n_1n_2$, so $n_3$ has no edge at all:
\[
A = \begin{pmatrix} 0&1&0\\ 1&0&0\\ 0&0&0 \end{pmatrix},\quad
A^2 = \begin{pmatrix} 1&0&0\\ 0&1&0\\ 0&0&0 \end{pmatrix},\quad
S = \begin{pmatrix} 1&1&0\\ 1&1&0\\ 0&0&0 \end{pmatrix}.
\]
Here $S_{13} = S_{23} = S_{31} = S_{32} = 0$: no walk of length $1$ or $2$ joins $n_3$ to another vertex, and indeed $n_3$ has no neighbour, so no walk of any length leaves it. By \Cref{thm:test} the graph is not connected. The positive entry $S_{11} = 1$ (the walk $n_1, n_2, n_1$) does not help reach $n_3$, which is why the test reads the off-diagonal entries.
\end{example}

\section{What the test does not give}\label{sec:limits}

The entries of $S$ count walks, not paths, so only their positivity is used: on $G_1$, $S_{22} = 2$ counts the walks $n_2, n_1, n_2$ and $n_2, n_3, n_2$, but only one path, $(n_2)$, joins $n_2$ to itself (\Cref{ex:board}, \Cref{ex:test1}). The test is also not the cheapest way to decide connectivity. It forms $N - 2$ matrix products, while a graph search from one vertex such as depth-first search runs in time $O(E)$ for $E$ edges, since it examines each edge at most twice \cite{wayne_cos226_undirected}. That the test is therefore the more expensive one is an interpretation: the source states the running time for depth-first search only, says nothing about breadth-first search or about this test, and no complexity analysis is made here.

Two sources state these results with the word ``paths'' where walks are meant. Paramadevan and Sotheeswaran's Theorem 1 says the entry of $A^l$ is the number of ``different paths'' of length $l$ \cite[Theorem 1, p.~16]{paramadevan2021adjacency} and the paper never defines ``path''; read as walks it is \Cref{thm:walks}, read with the paths of \Cref{def:path} it fails, since $(A^2)_{11} = 1$ on $G_1$ while no path of length $2$ joins $n_1$ to itself. Bai's proof sketch says $(I+A)^{n-1}$ ``aggregates all paths of length $\le n-1$'' \cite[p.~2]{bai_bimsa_adjacency}, while his Theorem 1 says that powers of $A$ count walks, not paths \cite[p.~1]{bai_bimsa_adjacency}; the entries of $(I+A)^{n-1}$ count walks too (\Cref{rem:binom}).

\section*{Notes on sources}

Questions 1 and 2 come from a lecture board photographed twice \cite{whiteboard1,whiteboard2}; question 3 comes from an assignment known only through a summary \cite{user_statements}, so the claim that the course's ``path'' means no repeated vertex rests on that summary and not on the assignment text, and was not checked against it. The two small marks above $A$ on the board are read as the exponent $k$, as the summary says. The board's drawing has three lines (1--2 with an arrowhead, 2--3, 1--3), while its matrix has two edges, is symmetric, and has a zero diagonal; the paper uses the matrix (\Cref{ex:board}). The board writes the graph as $(G, V)$ and the edge condition as ``$e_{ij}$ in $V$'', and defines connected with ``path'' undefined; the textbook writes $G = (V,E)$ and ``adjacent''.

The textbook contains no walk-counting theorem and no matrix test for connectivity on the pages viewed (printed pp.~11, 12, 34, 61), after a text search of the whole book for the adjacency-matrix phrase; page 7, a notation list, was not viewed, so this absence is not certain. Its authors are given by surname only, and publisher and year were not read. Grinberg and Bai use ``path'' as this paper does \cite{grinberg_graph_theory,bai_bimsa_adjacency}; the textbook's ``path'' is the trail. The scripts are in \texttt{scripts/}: \texttt{m1\_walks.py} and \texttt{m2\_connectivity\_test.py}, run with \texttt{python3 -I} from the paper folder, with their output files beside them.


\printbibliography


\end{document}
